\section{Introducción}

\subsection{Problema}

La delimitación del campo pulmonar precede a toda medición morfológica en una
radiografía de tórax. Su resultado depende del detector, de la exposición, del
procesamiento de la imagen y de los hallazgos que opacifican el parénquima. Si la
máscara obtenida se desplaza según la estrategia de preprocesamiento, los
descriptores calculados sobre ella heredan esa variabilidad y dejan de ser
comparables entre sujetos o entre sitios de adquisición, aunque la anatomía no haya
cambiado.

Las técnicas del curso permiten cuantificar ese efecto de manera controlada:
aplicar estrategias alternativas de realce, umbralización y morfología, medir su
concordancia con máscaras trazadas por personas y estimar cuánto se desplaza cada
descriptor respecto a su valor de referencia.

\subsection{Pregunta de investigación}

\begin{quote}
¿En qué medida las estrategias de realce de intensidad, umbralización y morfología
modifican la estabilidad de los descriptores morfológicos del campo pulmonar
(compacidad, relación de aspecto, fracción de área, simetría izquierda-derecha y
energía wavelet), respecto a los descriptores obtenidos sobre la máscara de
referencia, en radiografías de tórax de los conjuntos Montgomery y Shenzhen?
\end{quote}

\subsection{Objetivos}

\textbf{Objetivo general.} Evaluar el efecto de las estrategias de normalización,
realce, umbralización y morfología sobre la estabilidad de los descriptores
morfológicos del campo pulmonar, mediante la comparación cuantitativa con los
descriptores derivados de máscaras de referencia, para determinar qué combinación
de técnicas produce mediciones reproducibles entre sitios de adquisición.

\begin{table}[H]
\centering
\small
\caption{Objetivos específicos, fase y sección de este informe}
\begin{tabular}{@{}p{0.9cm}p{9.6cm}cc@{}}
\toprule
\textbf{OE} & \textbf{Objetivo} & \textbf{Fase} & \textbf{Sección} \\
\midrule
OE1 & Caracterizar la digitalización y adquisición de ambos conjuntos: resolución,
      cuantización, canales e intensidades & 1 & \ref{sec:caracterizacion} \\
OE2 & Comparar estrategias de normalización y reducción de ruido en dominio
      espacial y de frecuencia sobre una muestra controlada & 1--2 & \ref{sec:res-oe2} \\
OE3 & Comparar estrategias de segmentación del campo pulmonar contra las máscaras
      de referencia & 2 & \ref{sec:res-oe3} \\
OE4 & Cuantificar el desplazamiento de cada descriptor respecto a su valor de
      referencia & 2--3 & \ref{sec:res-oe4} \\
OE5 & Analizar las diferencias entre grupos con y sin hallazgos y documentar las
      limitaciones & 3 & \ref{sec:res-oe5} \\
\bottomrule
\end{tabular}
\end{table}

Además, el laboratorio de modelado supervisado agregó una pregunta predictiva:
si los descriptores de la segmentación automática permiten estimar la fracción
pulmonar de la máscara manual en pacientes no vistos (sección \ref{sec:res-modelo}).

\subsection{Alcance}

El proyecto caracteriza y compara; \textbf{no detecta ni diagnostica tuberculosis}.
La etiqueta clínica se usa solo para comparar grupos (OE5). No propone un método
superior al estado del arte, no evalúa observadores humanos y no generaliza sus
resultados a otras modalidades ni a otras proyecciones.

\subsection{Organización del trabajo}

El proyecto se desarrolló en tres fases a lo largo de quince semanas, y cada
laboratorio del curso extendió el mismo repositorio:

\begin{table}[H]
\centering
\small
\caption{Correspondencia entre los laboratorios y este informe}
\begin{tabular}{@{}p{5.2cm}p{6.4cm}c@{}}
\toprule
\textbf{Laboratorio} & \textbf{Producto principal} & \textbf{Sección} \\
\midrule
Formulación del proyecto & Pregunta, objetivos, datos, matriz y cronograma & 1 \\
Caracterización (Fase 1) & Manifiesto, calidad y partición & 2 \\
Análisis exploratorio y tabla de características & Nueve estrategias y tabla de
descriptores & 3 y 4 \\
Modelado supervisado & Regresión múltiple con protocolo congelado & 4 \\
Fase 3: OE2 a OE5 & Realce, frecuencia, wavelets, desplazamiento y grupos & 4 \\
Laboratorios preparatorios del curso & Entorno reproducible y adquisición & Anexo A \\
\bottomrule
\end{tabular}
\end{table}